%%% EDITÁVEL%%%
%numero de autores máximo 3
\newcounter{numerodeautores}\addtocounter{numerodeautores}{4}% numero 1, 2 ou 3 
\newcommand{\titulonotatecnica}{Título da Nota Técnica}
\newcommand{\autornotatecnica}{Autor1}
\newcommand{\filiacaoautordanota}{Laboratório Nacional de Astrofísica}
\newcommand{\emailautor}{autor1@lna.br} %necessario apenas para o primeiro autor
\newcommand{\numeronotatecnica}{LNA/NT/202X-z}
\newcommand{\datanotatecnica}{Mês/ano}
\newcommand{\segundoautor}{Autor2}
\newcommand{\terceiroautor}{Autor3}
\newcommand{\quartoautor}{Autor4}

%Revisores em ordem alfabética
\newcommand{\revisorum}{Revisor 1}
\newcommand{\revisordois}{Revisor 2}
\newcommand{\revisortres}{Revisor 3 -- nomes dos revisores em ordem alfabética}

% Resumo %
\newcommand{\resumonotatecnica}{Aqui vai um resumo do trabalho a ser apresentado ao longo da nota técnica}

% Abstract %
\newcommand{\abstractnotatecnica}{\textit{Include here the abstract of your work}}

% Palavras-chave/keywords
\newcommand{\keywordsnotatecnica}{coloque aqui as palavras-chave para futura indexação. Exemplo: espectrógrafo; rede echelle; projeto óptico / \textit{spectrograph; echelle grating; optical design}}